\section{Experiments}
\label{sec:experiments}

\subsection{Datasets}

\paragraph{CIFAR-100}~\cite{krizhevsky2009learning} contains 60,000 color images of size $32 \times 32$ pixels across 100 classes, with 500 training and 100 test images per class. Images are resized to $224 \times 224$ before feature extraction to match the input resolution expected by ImageNet-pretrained networks.

\paragraph{Tiny ImageNet}~\cite{le2015tiny} is a subset of ImageNet with 200 classes, each containing 500 training images and 50 validation images at $64 \times 64$ pixels. Images are similarly resized to $224 \times 224$.

The two datasets provide complementary evaluation conditions: CIFAR-100 tests the low-resolution regime with $C - 1 = 99$ LDA components, while Tiny ImageNet provides more classes ($C - 1 = 199$), more training data, and higher native resolution.

\subsection{Backbone Architectures}

We extract features from four pretrained ImageNet classifiers spanning a range of architectures and feature dimensions:

\begin{itemize}[leftmargin=*]
    \item \textbf{ResNet-18}~\cite{he2016resnet} ($D = 512$): a compact residual network, widely used as a baseline.
    \item \textbf{ResNet-50}~\cite{he2016resnet} ($D = 2048$): a deeper variant with 4$\times$ higher feature dimensionality.
    \item \textbf{MobileNetV3-Small}~\cite{howard2019mobilenetv3} ($D = 576$): designed for mobile deployment; lightweight architecture.
    \item \textbf{EfficientNet-B0}~\cite{tan2019efficientnet} ($D = 1280$): a modern compound-scaled network balancing depth, width, and resolution.
\end{itemize}

All backbones use ImageNet-pretrained weights from torchvision~\cite{pytorch}. Features are extracted from the layer immediately before the classification head (global average pooling output), with all backbone parameters frozen. Features are extracted once and cached.

\subsection{Experimental Protocol}

For each backbone--dataset pair:
\begin{enumerate}[leftmargin=*]
    \item Extract features from all training and test images.
    \item Fit the dimensionality reduction method on training features only.
    \item Transform both training and test features.
    \item Standardize features (zero mean, unit variance; fit on training set).
    \item Train an $\ell_2$-regularized logistic regression classifier (LBFGS solver, $\text{max\_iter}=5000$, $C=1.0$).
    \item Evaluate top-1 accuracy on the test set.
\end{enumerate}

All reduction methods project to $d = C - 1$ dimensions (99 for CIFAR-100, 199 for Tiny ImageNet), except RDA which produces $(C-1) + k$ dimensions. Timing is measured using wall-clock time with \texttt{time.perf\_counter()} and includes both the reduction step and classifier training.

For statistical significance, we repeat each configuration across 5 random seeds (affecting classifier initialization and any stochastic components) and report mean accuracy with standard deviations.

\subsection{Main Results}

\subsubsection{CIFAR-100}

Table~\ref{tab:cifar100} reports accuracy and total wall-clock time for all ten methods across four backbones on CIFAR-100.

\begin{table*}[t]
\centering
\caption{Classification accuracy (\%, mean$\pm$std over 5 seeds) and total time (seconds) on \textbf{CIFAR-100} across four backbones. The best result per backbone is \textbf{bolded}; the second best is \underline{underlined}. All reduction methods project to $d = 99$ dimensions ($C - 1$) except RDA variants which use 99$+k$. Methods marked with $^\dagger$ are single-seed results.}
\label{tab:cifar100}
\small
\begin{tabular}{@{}l l cc cc cc cc@{}}
\toprule
& & \multicolumn{2}{c}{\textbf{ResNet-18} (512D)} & \multicolumn{2}{c}{\textbf{ResNet-50} (2048D)} & \multicolumn{2}{c}{\textbf{MobileNetV3} (576D)} & \multicolumn{2}{c}{\textbf{EfficientNet} (1280D)} \\
\cmidrule(lr){3-4} \cmidrule(lr){5-6} \cmidrule(lr){7-8} \cmidrule(lr){9-10}
Cat. & Method & Acc. & Time & Acc. & Time & Acc. & Time & Acc. & Time \\
\midrule
\multirow{2}{*}{Control}
& Full       & 62.85{\tiny$\pm$.02} & 56   & 71.98{\tiny$\pm$.00} & 48   & 65.58{\tiny$\pm$.07} & 69   & 71.41{\tiny$\pm$.04} & 49  \\
& PCA        & 65.06{\tiny$\pm$.01} & 11   & 69.18{\tiny$\pm$.12} & 15   & 64.61{\tiny$\pm$.01} & 11   & 69.77{\tiny$\pm$.13} & 14  \\
\midrule
\multirow{2}{*}{Classical}
& LDA        & 66.97{\tiny$\pm$.00} & 11   & 72.24{\tiny$\pm$.05} & 23   & 68.47{\tiny$\pm$.03} & 13   & 72.30{\tiny$\pm$.04} & 16  \\
& PCA+LDA$^\dagger$    & 66.63 & 10   & 70.68 & 19   & 67.97 & 12   & 70.80 & 17  \\
\midrule
\multirow{3}{*}{Academic}
& R-LDA      & 66.68{\tiny$\pm$.02} & 13   & 71.89{\tiny$\pm$.06} & 43   & 67.86{\tiny$\pm$.04} & 15   & 71.95{\tiny$\pm$.03} & 24  \\
& LFDA       & 65.25{\tiny$\pm$.00} & 15   & 68.20{\tiny$\pm$.09} & 22   & 64.90{\tiny$\pm$.03} & 19   & 69.33{\tiny$\pm$.11} & 20  \\
& NCA$^\dagger$        & 64.93 & 258  & 69.32 & 267  & 65.71 & 283  & 69.95 & 269 \\
\midrule
\multirow{3}{*}{Ours}
& RDA        & 66.75{\tiny$\pm$.07} & 25   & \underline{72.67{\tiny$\pm$.05}} & 44   & \underline{68.72{\tiny$\pm$.04}} & 29   & 72.35{\tiny$\pm$.03} & 35  \\
& DSB        & \textbf{67.27{\tiny$\pm$.00}} & 25   & 72.62{\tiny$\pm$.06} & 58   & \textbf{68.88{\tiny$\pm$.04}} & 33   & \textbf{72.55{\tiny$\pm$.00}} & 42  \\
& RDA+SMD$^\dagger$    & 66.74 & 37   & \textbf{73.06} & 182  & 68.73 & 51   & \underline{72.40} & 76  \\
\bottomrule
\end{tabular}
\end{table*}

Several patterns emerge. First, LDA improves over full features on all four backbones: by +4.12\% on ResNet-18, +0.26\% on ResNet-50, +2.89\% on MobileNetV3, and +0.89\% on EfficientNet. The gains are most pronounced on compact backbones (ResNet-18, MobileNetV3), where the full-feature classifier struggles most and LDA's discriminant projection provides the greatest lift.

Second, LDA strictly dominates PCA: the supervisory signal consistently translates into higher accuracy. The gap ranges from +1.91\% (ResNet-18) to +3.86\% (MobileNetV3). This confirms that label information provides value beyond variance-based feature selection.

Third, more complex academic methods provide no advantage over plain LDA. R-LDA slightly underperforms LDA in three of four cases---the Ledoit--Wolf shrinkage adds computation without benefit because scikit-learn's SVD solver already handles conditioning implicitly. LFDA performs poorly across the board, 1.7--4.0 points below LDA, suggesting that the multimodal assumption does not hold for frozen CNN features, which tend to form unimodal, approximately Gaussian class clusters. NCA is 10--25$\times$ slower than LDA while delivering lower accuracy.

Fourth, our extensions (RDA, DSB) can improve on LDA, but the gains are modest: DSB adds 0.30\% on ResNet-18 at 2.3$\times$ the cost, and RDA+SMD adds 0.82\% on ResNet-50 at 8$\times$ the cost. Whether these marginal improvements justify the extra computation depends on the application.

\subsubsection{Tiny ImageNet}

Table~\ref{tab:tinyimagenet} shows results on Tiny ImageNet, which has 200 classes and twice the training data.

\begin{table*}[t]
\centering
\caption{Classification accuracy (\%, mean$\pm$std over 5 seeds) and total time (seconds) on \textbf{Tiny ImageNet} across four backbones. Formatting follows Table~\ref{tab:cifar100}. All reduction methods project to $d = 199$ dimensions.}
\label{tab:tinyimagenet}
\small
\begin{tabular}{@{}l l cc cc cc cc@{}}
\toprule
& & \multicolumn{2}{c}{\textbf{ResNet-18} (512D)} & \multicolumn{2}{c}{\textbf{ResNet-50} (2048D)} & \multicolumn{2}{c}{\textbf{MobileNetV3} (576D)} & \multicolumn{2}{c}{\textbf{EfficientNet} (1280D)} \\
\cmidrule(lr){3-4} \cmidrule(lr){5-6} \cmidrule(lr){7-8} \cmidrule(lr){9-10}
Cat. & Method & Acc. & Time & Acc. & Time & Acc. & Time & Acc. & Time \\
\midrule
\multirow{2}{*}{Control}
& Full       & 59.84{\tiny$\pm$.04} & 122  & 74.29{\tiny$\pm$.00} & 84   & 59.24{\tiny$\pm$.06} & 150  & 71.79{\tiny$\pm$.00} & 68  \\
& PCA        & 64.46{\tiny$\pm$.02} & 30   & 73.35{\tiny$\pm$.08} & 52   & 61.95{\tiny$\pm$.04} & 35   & 71.99{\tiny$\pm$.17} & 46  \\
\midrule
\multirow{2}{*}{Classical}
& LDA        & 64.42{\tiny$\pm$.06} & 41   & 75.01{\tiny$\pm$.01} & 80   & 63.35{\tiny$\pm$.01} & 46   & 72.34{\tiny$\pm$.06} & 53  \\
& PCA+LDA$^\dagger$    & 64.29 & 42   & 74.03 & 77   & 63.30 & 50   & 72.01 & 61  \\
\midrule
\multirow{3}{*}{Academic}
& R-LDA      & 64.74{\tiny$\pm$.05} & 50   & \underline{75.37{\tiny$\pm$.09}} & 118  & 62.94{\tiny$\pm$.06} & 48   & \underline{72.58{\tiny$\pm$.07}} & 58  \\
& LFDA       & 64.15{\tiny$\pm$.03} & 49   & 73.54{\tiny$\pm$.12} & 73   & 62.57{\tiny$\pm$.00} & 50   & 71.92{\tiny$\pm$.04} & 64  \\
& NCA$^\dagger$        & 64.27 & 278  & 73.13 & 300  & 62.03 & 295  & 71.91 & 328 \\
\midrule
\multirow{3}{*}{Ours}
& RDA        & 64.52{\tiny$\pm$.03} & 44   & 75.04{\tiny$\pm$.06} & 95   & \underline{63.46{\tiny$\pm$.02}} & 49   & 71.87{\tiny$\pm$.06} & 71  \\
& DSB        & \textbf{64.70{\tiny$\pm$.00}} & 88   & 75.24{\tiny$\pm$.06} & 218  & \textbf{63.73{\tiny$\pm$.00}} & 87   & 72.14{\tiny$\pm$.06} & 147 \\
& RDA+SMD$^\dagger$    & 64.47 & 69   & 74.96 & 621  & 63.41 & 61   & 72.10 & 172 \\
\bottomrule
\end{tabular}
\end{table*}

The Tiny ImageNet results reinforce the patterns observed on CIFAR-100, with some interesting differences. LDA again beats full features on all four backbones, with the gains on lower-capacity architectures being particularly striking: +4.58 points on ResNet-18 and +4.11 on MobileNetV3. On ResNet-50, LDA achieves 75.01\% versus 74.29\% for full features---a more modest but still consistent improvement.

A notable difference is that PCA performs relatively better on Tiny ImageNet, particularly on ResNet-18 where it slightly exceeds LDA (64.46 vs. 64.42). With 200 classes and 100,000 training samples, the variance structure captured by PCA is richer, partially compensating for the lack of supervisory information.

R-LDA shows stronger performance on Tiny ImageNet than on CIFAR-100, achieving the highest accuracy on ResNet-50 (75.37\%) and EfficientNet (72.58\%). With 200 classes and higher-dimensional features, the scatter matrix conditioning becomes more of a concern, and regularization provides a genuine benefit.

RDA+SMD, which was the top method on ResNet-50/CIFAR-100, does not generalize its advantage to Tiny ImageNet. Its time cost on ResNet-50 is extreme (621s vs. 80s for LDA), making it impractical at scale.

\subsection{Summary of Main Results}

Across all eight backbone--dataset configurations:
\begin{itemize}[leftmargin=*]
    \item \textbf{LDA beats Full features in 8/8 cases} (+0.26\% to +4.58\%), while reducing dimensionality by 61--95\%.
    \item \textbf{LDA beats PCA in 7/8 cases.} The exception is ResNet-18 on Tiny ImageNet, where PCA edges LDA by 0.04\%.
    \item \textbf{LDA beats or matches LFDA in 8/8 cases} and \textbf{beats NCA in 8/8 cases}, while being 3--25$\times$ faster.
    \item \textbf{DSB achieves the highest accuracy in 4/8 cases} (both MobileNetV3 configs, ResNet-18/CIFAR-100, and EfficientNet/CIFAR-100), at 2--3$\times$ LDA's cost.
    \item \textbf{R-LDA is the strongest academic variant}, competitive with LDA on CIFAR-100 and slightly better on some Tiny ImageNet configs.
\end{itemize}

Fig.~\ref{fig:accuracy_gain} visualizes the accuracy gain of PCA, LDA, and DSB over Full features across all eight configurations, highlighting LDA's consistent improvement.

\begin{figure}[t]
    \centering
    \includegraphics[width=\columnwidth]{fig1_accuracy_gain.pdf}
    \caption{Accuracy improvement over full features for PCA, LDA, and DSB across all eight backbone--dataset configurations. LDA consistently outperforms full features (all bars positive), with the largest gains on compact backbones.}
    \label{fig:accuracy_gain}
\end{figure}
